% Judul Capstone
\judul{Prototipe Sistem Pemantauan Kualitas Udara Kabin Kendaraan Berbasis Sensor CO dan CO\textsubscript{2} dengan Peringatan dan Aktuasi Otomatis}
\title{Prototype of an In-Vehicle Cabin Air Quality Monitoring System Based on CO and CO\textsubscript{2} Sensors with Automatic \mbox{Warning} and Actuation Mechanism}

% Jenis Dokumen
%% Jenis Dokumen : IMPLEMENTASI SOLUSI

% Kode Dokumen
%% Kode Dokumen : C-501

% Nomor Dokumen (Kode Tim/Kelompok Capstone)
\NoDok{B\_09}

% Nomor Revisi
\NoRev{00}

% Tanggal Penerbitan Dokumen
%% Otomatis terisi tanggal ketika file LaTeX ini di-compile

% Data Mahasiswa Capstone
%% Ketua kelompok : \MHS{<Nama Lengkap>}{<NIM>}{<Prodi>}{<Alamat Email>}
\MHS{Aurellya Ratna Dewanti}{23/517176/TK/56870}
	  {Teknologi Informasi}{aurellyaratnadewanti@mail.ugm.ac.id}

%% Anggota kelompok : \Anggota{<Nomor 1-3>}{<Nama Lengkap>}{<NIM>}{<Prodi>}{<Alamat Email>}
\Anggota{1}{Joecelyn Aurora Majesty}{23/514716/TK/56510}
	  {Teknologi Informasi}{joecelynauroramajesty@mail.ugm.ac.id}
\Anggota{2}{Maritza Vania Adelia}{23/517643/TK/56944}
	  {Teknologi Informasi}{maritzavaniaadelia@mail.ugm.ac.id}
\Anggota{3}{Muhammad Akbar Rabbani}{23/515815/TK/56739}
	  {Teknik Elektro}{muhammadakbarrabbani@mail.ugm.ac.id}

% Teks "PROGRAM STUDI" pada halaman judul (mahasiswa berasal dari dua program studi)
\ProdiJudul{Teknologi Informasi\\dan Teknik Elektro}

% Dosen Pembimbing
%% Format : {<Nama Lengkap>}{<NIP/NIU>}
\DPA{Dr. Ir. I Wayan Mustika, S.T., M.Eng.}{198109212014041001}

% Tempat Pelaksanaan
%% Format : {<Nama Laboratorium> \newline <Nama Departemen> \newline <Nama Fakultas>}
\Tempat{Departemen Teknik Elektro dan Teknologi Informasi \newline Fakultas Teknik}
